\paragraph{固定分辨率全息完备性、屏幕相对同调与 Artin 通道商擦除}
在固定分辨率 dyadic 体--界全息、部分屏幕观测的相对同调核、Zeckendorf 折叠的确定化歧义壳以及非阿贝尔 Peter--Weyl/Artin 分块都已建立之后，还可进一步抽取出下列七条后果。它们共同表明：困难并不首先来自 bulk 无法在有限层面被完整编码；真正拒绝统一有限压缩的，是部分屏幕留下的相对同调赤字，以及群商压缩所擦除的非阿贝尔 Artin 通道。

\begin{theorem}[固定分辨率 Stokes--Gödel 全息完备性的强形式]\label{thm:xi-time-part9x-fixed-resolution-stokes-godel-strong-completeness}
固定 \(m,n\ge 1\)，令 \(\mathcal U_{m,n}\) 为 \(I^n=[0,1]^n\) 上全部分辨率 \(m\) 的 dyadic polycube，\(\mathcal Q_m\) 为对应的 \(m\)-级 dyadic 立方复形。对每个 \(U\in\mathcal U_{m,n}\)，记 \(u_U\in C_n(\mathcal Q_m;\FF_2)\) 为其占据链，并令
\[
\Phi(U):=(G\circ\partial_n)(u_U),
\]
其中 \(G\) 为任意在像上可计算可逆的平方自由 boundary Gödel 整数化。则存在常数 \(C>0\)（仅依赖于所选通用前缀机与编码方案）使得对一切 \(U\in\mathcal U_{m,n}\)，都有
\[
\bigl|K(U\mid m,n)-K(\Phi(U)\mid m,n)\bigr|\le C,
\qquad
K(U\mid \Phi(U),m,n)\le C.
\]
此外，\(\Phi(U)\) 唯一确定 \(U\)，并且唯一确定有限矩盒
\[
\{\mu_\alpha(U):0\le \alpha_i\le 2^m-1\}.
\]
因而在固定分辨率上，boundary Gödel code 是复杂度守恒且信息完备的全息坐标。
\end{theorem}
\begin{proof}
在给定 \((m,n)\) 时，占据链 \(u_U\) 与 polycube \(U\) 之间存在可计算双向转换，故
\[
K(U\mid m,n)=K(u_U\mid m,n)+O(1).
\]
由定理 \ref{thm:spg-stokes-godel-algorithmic-holographic-completeness}，
\[
\bigl|K(u_U\mid m,n)-K((G\circ\partial_n)(u_U)\mid m,n)\bigr|\le O(1),
\qquad
K(u_U\mid \Phi(U),m,n)\le O(1).
\]
把 \(U\leftrightarrow u_U\) 的编码常数吸收进同一 \(O(1)\) 项，即得复杂度结论。

另一方面，推论 \ref{cor:spg-boundary-godel-finite-moment-completeness} 已表明 \(\Phi(U)\) 唯一决定 \(U\)。更具体地，由定理 \ref{thm:spg-boundary-godel-moment-readout}，\(\Phi(U)\) 决定全部单项式矩；再由定理 \ref{thm:spg-dyadic-finite-moment-completeness}，其中有限矩盒已足以唯一恢复 \(U\)。证毕。
\end{proof}

\begin{theorem}[部分屏幕观测的相对同调赤字闭式]\label{thm:xi-time-part9x-partial-screen-relative-homology-deficit-closed}
沿用定理 \ref{thm:spg-partial-boundary-screen-kernel-relative-homology} 的记号。对任意屏幕
\[
S\subset \overrightarrow{\mathcal F_{m}^{n-1}},
\]
其部分边界观测
\[
f_S:=\Pi_S\circ\partial_n:\ C_n(\mathcal Q_m;\ZZ)\longrightarrow \ZZ^S
\]
满足自然同构
\[
\ker f_S\cong H_n(\mathcal Q_m,A_{\neg S};\ZZ)\cong \ZZ^{\,c(\Gamma_S)-1},
\]
其中 \(\Gamma_S\) 为屏幕图，\(c(\Gamma_S)\) 为其连通分量数。特别地，
\[
\rank_{\ZZ}\ker f_S=c(\Gamma_S)-1.
\]
\end{theorem}
\begin{proof}
前一同构由定理 \ref{thm:spg-partial-boundary-screen-kernel-relative-homology} 给出，后一同构由定理 \ref{thm:spg-screen-kernel-connected-components} 给出。复合即得结论。证毕。
\end{proof}
\leanverified{paper\_xi\_time\_part9x\_partial\_screen\_relative\_homology\_deficit\_closed}

\begin{corollary}[屏幕补全成本的面积律 sharp 化]\label{cor:xi-time-part9x-screen-completion-cost-sharp-area-law}
沿用定理 \ref{thm:xi-time-part9x-partial-screen-relative-homology-deficit-closed} 的记号，定义最小审计补全成本
\[
\mathsf{Cost}(S):=\min_{(R,r_S)}\cdim(R)=\rank_{\ZZ}\ker f_S.
\]
则
\[
\mathsf{Cost}(S)=c(\Gamma_S)-1.
\]
特别地，对轴向屏幕 \(S^{(1)}\)（命题 \ref{prop:spg-axial-screen-area-law-audit-cost} 的记号），有精确公式
\[
\mathsf{Cost}(S^{(1)})=2^{m(n-1)}.
\]
\end{corollary}
\begin{proof}
由推论 \ref{cor:spg-screen-injectivity-and-audit-cost-components}，
\[
\mathsf{Cost}(S)=\rank_{\ZZ}\ker f_S=c(\Gamma_S)-1.
\]
对轴向屏幕 \(S^{(1)}\)，命题 \ref{prop:spg-axial-screen-area-law-audit-cost} 给出
\[
c(\Gamma_{S^{(1)}})=2^{m(n-1)}+1.
\]
代入第一式即可。证毕。
\end{proof}

\begin{theorem}[瞬态壳与周期谱的严格正交]\label{thm:xi-time-part9x-transient-shell-periodic-spectrum-orthogonality}
对 Zeckendorf 折叠的右确定化表示 \(H_m^{\det}\)（\(m\ge 3\)），按“歧义态在前、单点态在后”的顺序排列状态后，其邻接矩阵可写为
\[
A_m^{\det}
=
\begin{pmatrix}
N_m & R_m\\
0 & A_m^{\mathrm{ess}}
\end{pmatrix},
\]
其中 \(N_m\) 为 ambiguity shell 块，\(A_m^{\mathrm{ess}}\) 为 singleton skeleton 块。则成立：
\[
N_m^{\,m}=0,
\qquad
\det(I-zA_m^{\det})=\det(I-zA_m^{\mathrm{ess}}),
\qquad
D_{\mathrm{sync}}(m)=m-1.
\]
因此歧义壳既不承载任何周期轨道，也不改变表示图层面的 \(\zeta\)-因子；它只以有限同步延迟进入时间复杂度。

更一般地，在一般 sofic 口径下，若歧义壳子系统 \(Y_m^{\mathrm{amb}}\neq\varnothing\) 且其最大本质分量的周期为 \(p\ge 1\)，则存在常数族
\[
\{c_r(m)\}_{r=0}^{p-1}\subset(0,\infty)
\]
使得
\[
\nu_m(T_m>n)=c_{n\bmod p}(m)e^{-\varepsilon_m n}(1+o(1)),
\qquad
\varepsilon_m:=\log 2-\log\rho_{\mathrm{amb}}(m).
\]
\end{theorem}
\begin{proof}
前两式由定理 \ref{thm:Ym-ambiguity-shell-zeta-factor-trivial} 直接给出，其中 \(N_m\) 幂零且 \(\det(I-zN_m)=1\)。同步延迟闭式由推论 \ref{cor:Ym-sync-delay} 给出。最后一式正是定理 \ref{thm:Ym-sync-tail-pf-sharp} 的结论。故周期/Fredholm 数据与同步预算在结构上正交：前者只感知单点骨架，后者则由歧义壳的瞬态或低熵尾主导。证毕。
\end{proof}

\begin{theorem}[非阿贝尔 Artin 通道在商压缩下的 determinant--energy 双擦除]\label{thm:xi-time-part9x-nonabelian-artin-channel-quotient-erasure}
沿用命题 \ref{prop:kernel-peter-weyl-block-diagonalization} 与命题 \ref{prop:kernel-artin-factorization} 的记号。设 \(G\) 为有限群，\(K\in\CC[G]^{V\times V}\) 为群标记核，\(\mathcal K\in\End(\CC^V\otimes\CC[G])\) 为其右卷积实现，并记
\[
\mathcal K\cong \bigoplus_{\rho\in\widehat G}\bigl(K_\rho\otimes I_{d_\rho}\bigr),
\qquad
\det(I-z\mathcal K)=\prod_{\rho\in\widehat G}\det(I-zK_\rho)^{d_\rho},
\]
其中 \(d_\rho:=\dim\rho\)。对任意满同态
\[
\pi:G\twoheadrightarrow H,
\]
令 \(\pi_*K\in\CC[H]^{V\times V}\) 为逐矩阵项推前的商核，\(\mathcal K_H\) 为其右卷积实现，并定义
\[
\mathrm{Inf}_\pi(\widehat H):=\{\sigma\circ\pi:\sigma\in\widehat H\}\subseteq \widehat G.
\]
则对每个 \(\sigma\in\widehat H\) 都有
\[
(\pi_*K)_\sigma=K_{\sigma\circ\pi},
\]
从而
\[
\det(I-z\mathcal K_H)=\prod_{\sigma\in\widehat H}\det(I-zK_{\sigma\circ\pi})^{d_\sigma},
\]
其中 \(d_\sigma:=\dim\sigma\)。特别地，一切不属于 \(\mathrm{Inf}_\pi(\widehat H)\) 的 irreducible Artin channel 都从商核的 determinant 因子中完全消失。

此外，对任意零均值类函数 \(\varepsilon\in\mathrm{Cl}(G)\)，其纤维平均推前满足精确能量损失公式
\[
\|\varepsilon\|_{L^2(G)}^2-\|\pi_*\varepsilon\|_{L^2(H)}^2
=
\frac1{|G|^2}
\sum_{\rho\in\widehat G\setminus\mathrm{Inf}_\pi(\widehat H)}
\bigl|B_G(\rho)\bigr|^2.
\]
因此群商压缩会把所有非 inflation 通道同时从 determinant 数据与非阿贝尔能量账本中不可逆擦除。
\end{theorem}
\begin{proof}
对任意群代数元素
\[
a=\sum_{g\in G}a_g\,g\in\CC[G]
\]
及任意 \(\sigma\in\widehat H\)，逐项推前定义给出
\[
\sum_{h\in H}(\pi_*a)_h\,\sigma(h)
=
\sum_{h\in H}\left(\sum_{\pi(g)=h}a_g\right)\sigma(h)
=
\sum_{g\in G}a_g\,\sigma(\pi(g))
=
\sum_{g\in G}a_g\,(\sigma\circ\pi)(g).
\]
逐矩阵项应用此恒等式，即得
\[
(\pi_*K)_\sigma=K_{\sigma\circ\pi}.
\]
再对 \(\pi_*K\) 应用命题 \ref{prop:kernel-artin-factorization}，便得到所示 determinant 分解；其中只出现 inflated 通道，因此所有 \(\rho\notin\mathrm{Inf}_\pi(\widehat H)\) 的因子都在商核侧消失。

能量损失公式则正是定理 \ref{thm:xi-time-part64-nonabelian-quotient-energy-loss-exact}。最后，由 Peter--Weyl 直和分解可知：若固定全部 inflated 通道而改变某个 \(\rho\notin\mathrm{Inf}_\pi(\widehat H)\) 的分块 \(K_\rho\)，则 \(\mathcal K_H\) 保持不变，而全 determinant 因子化与上式右端一般随之改变。故这些通道无法由商压缩后的数据重建。证毕。
\end{proof}

\begin{corollary}[阿贝尔化证书的严格不足性]\label{cor:xi-time-part9x-abelianized-certificate-strict-insufficiency}
令
\[
q:G\twoheadrightarrow G^{\mathrm{ab}}:=G/[G,G]
\]
为阿贝尔化商，并记 \(\mathcal K\) 为 \(K\) 的右卷积实现。若存在某个 \(\rho\in\widehat G\) 满足
\[
\dim\rho>1,
\qquad
\det(I-zK_\rho)\not\equiv 1,
\]
则任何仅依赖于 abelianized labels 的 Dirichlet/Artin 证书都不可能恢复完整 determinant
\[
\det(I-z\mathcal K),
\]
也不可能恢复由全部 Artin 因子共同编码的 twisted primitive-orbit law。
\end{corollary}
\begin{proof}
对阿贝尔群 \(G^{\mathrm{ab}}\)，全部不可约表示均为一维表示。因此 \(\mathrm{Inf}_q(\widehat{G^{\mathrm{ab}}})\) 恰好是 \(G\) 的全部一维不可约表示。由定理 \ref{thm:xi-time-part9x-nonabelian-artin-channel-quotient-erasure}，所有满足 \(\dim\rho>1\) 的通道都会从 \(q_*K\) 的 determinant 因子化中消失。若其中某个通道的 Artin 因子 \(\det(I-zK_\rho)\) 非平凡，则它不可能由 abelianized system 重建；故完整 determinant 与相应的 twisted primitive-orbit 统计也都无法恢复。证毕。
\end{proof}

\begin{conjecture}[classical RH 证书的 screen--Artin 双障碍]\label{conj:xi-time-part9x-rh-screen-artin-double-obstruction}
在本文的 solenoidal--Mellin--Fredholm 语义中，任何可审计的 classical RH 终端证书若要在共尾分辨率上保持完备，必须同时保留两类对象：
\[
\mathfrak H_{\mathrm{scr}}
:=
\{\,H_n(\mathcal Q_m,A_{\neg S_m};\ZZ)\,\}_{m\ge 1},
\qquad
\mathfrak H_{\mathrm{Art}}
:=
\{\,K_{\rho,m}:\rho\in\widehat G_m\setminus\mathrm{Inf}_{\pi_m}(\widehat H_m)\,\}_{m\ge 1}.
\]
其中 \(K_{\rho,m}\) 表示第 \(m\) 级群标记核的 \(\rho\)-通道块，而 \(\pi_m:G_m\twoheadrightarrow H_m\) 表示相应的商压缩。
更具体地说：
\begin{enumerate}
  \item 任何试图把 bulk 证书压缩到一族部分屏幕 \(S_m\) 且在共尾分辨率上保持统一有界补全成本的方案，只要 \(\rank H_n(\mathcal Q_m,A_{\neg S_m};\ZZ)\) 不最终消失，就必在某一级失去相对同调核信息；
  \item 任何试图把群标签统一压缩到固定有限商或纯阿贝尔化的方案，只要仍存在非 inflation 的活跃 Artin 通道，就必丢失相应的 determinant 因子与通道能量；
  \item 因而，同时施行这两类 finite-horizon 压缩的正规化，不应能够闭合 classical RH 所需的完整终端证书链。
\end{enumerate}
\end{conjecture}

\noindent
上述结果表明，本文对象中的困难并不首先来自“无限维对象不可构造”。在固定分辨率上，bulk 甚至可以由单个 boundary Gödel 整数复杂度守恒地编码并唯一恢复。真正不可统一有限化的，是部分 screen 观测留下的相对同调赤字，以及群商压缩所抹去的非阿贝尔 Artin determinant 因子。因而，更合适的终端障碍语言应当同时记录 screen completion cost 与 nonabelian channel loss；在此意义下，screen--Artin obstruction theory 比单纯的“有限见证不足”命题更接近本文对象的真实不可压缩核心。

\endinput
